\section{Discussion and Conclusion}
\label{sec:conclusion}

We asked whether adapting the CFR regret update to the local dynamics of each information set
improves convergence. Our results suggest that it does, but that \emph{what} is adapted
matters. Adapting only the discount on past regret (\ours{}) helped relative to \cfrp{} with the
same averaging, and its local version consistently beat its global version, which supports
the idea that information sets differ in how much past regret they should keep. However, its
gains were modest, depended on settings selected on the evaluation games, and did not
transfer to all held-out games. Letting each information set choose among existing update rules
(MIX-CFR) was much more effective. Without any tuning, it was the best or close to the best
method on all five games, while every fixed rule was far behind on at least one, and on two
games it outperformed all of the rules it combines.

\textbf{Limitations.} Our games are small enough for exact exploitability, and the runs are
limited to 1000--2000 iterations, so it remains open whether the advantage persists in much
larger games, in longer runs, or with sampling-based CFR. The local meta learner was clearly
better than the global one on Kuhn and Leduc poker but not on every game, so the benefit of
local selection also depends on the game. Our convergence guarantee covers the uniform average
strategy, not the $t^2$-weighted average used in the experiments. Finally, our timings come
from a Python implementation run in parallel and are approximate.

\textbf{Future work.} Natural extensions are adding more experts, such as PDCFR$^+$ or
APCFR$^+$, evaluating on larger games with a C++ implementation, and analyzing which features
of an information set determine the expert it selects.
